\expandafter\ifx\csname CAPETINCLUDEMODE\endcsname\relax
\documentclass[11pt,a4paper]{article}
% preamble.tex — packages et commandes communes a toutes les fiches CAPET SII IE
\input{glyphtounicode}
\pdfgentounicode=1
\usepackage{cmap}
\usepackage[T1]{fontenc}
\usepackage[utf8]{inputenc}
\usepackage[french]{babel}
\usepackage[a4paper, top=1.6cm, bottom=1.6cm, left=1.4cm, right=1.4cm]{geometry}
\setlength{\headheight}{14.5pt}
\usepackage{amsmath, amssymb, amsfonts}
\usepackage{siunitx}
\sisetup{locale=FR, output-decimal-marker={,}, inter-unit-product=\ensuremath{{}\cdot{}}}
\usepackage{array, booktabs, multirow, tabularx}
\usepackage{xcolor}
\usepackage[most, breakable]{tcolorbox}
\usepackage{multicol}
\usepackage{enumitem}
\usepackage{mathtools}
\usepackage{lmodern}
\usepackage{microtype}
\usepackage{fancyhdr}
\usepackage{tikz}
\usetikzlibrary{positioning,arrows.meta}
\usepackage{lastpage}

% ── Couleurs par domaine ──────────────────────────────────────────────
\definecolor{colElec}{RGB}{26, 82, 148}       % Électrotechnique — bleu foncé
\definecolor{colElecLight}{RGB}{210, 225, 245} % fond bleu clair
\definecolor{colAuto}{RGB}{34, 120, 66}        % Automatique — vert
\definecolor{colMeca}{RGB}{185, 90, 20}        % Mécanique — orange
\definecolor{colTherm}{RGB}{160, 40, 40}       % Thermique — rouge
\definecolor{colQT}{RGB}{232, 244, 233}        % fond question type (vert très léger)
\definecolor{colJury}{RGB}{255, 243, 205}      % fond jury (jaune ambre)
\definecolor{colJuryFrame}{RGB}{200, 130, 0}   % bordure jury
\definecolor{colGrayLight}{RGB}{245, 245, 247} % fond gris neutre
\definecolor{colSeparator}{RGB}{180, 180, 180} % gris séparateur

% ── Style tcolorbox de section principale ────────────────────────────
\tcbset{
  ficheSection/.style={
    breakable,
    enhanced,
    colback=#1!7!white,
    colframe=#1,
    colbacktitle=#1,
    coltitle=white,
    fonttitle=\bfseries\normalsize,
    left=4pt, right=4pt, top=3pt, bottom=4pt,
    before skip=6pt, after skip=6pt,
  }
}

% ── Style encadré question type ──────────────────────────────────────
\tcbset{
  questiontype/.style={
    breakable,
    enhanced,
    colback=colQT,
    colframe=colElec!55!black,
    colbacktitle=colElec!55!black,
    coltitle=white,
    fonttitle=\bfseries\small,
    left=4pt, right=4pt, top=3pt, bottom=3pt,
    before skip=5pt, after skip=5pt,
  }
}

% ── Style encadré jury ───────────────────────────────────────────────
\tcbset{
  jury/.style={
    breakable,
    enhanced,
    colback=colJury,
    colframe=colJuryFrame,
    fonttitle=\bfseries\normalsize,
    left=4pt, right=4pt, top=3pt, bottom=3pt,
  }
}

% ── Mise en page ─────────────────────────────────────────────────────
\pagestyle{fancy}
\fancyhf{}
\renewcommand{\headrulewidth}{0.4pt}
\renewcommand{\footrulewidth}{0.3pt}

\setlist[itemize]{leftmargin=1.2em, itemsep=1pt, topsep=2pt, parsep=0pt}
\setlist[enumerate]{leftmargin=1.2em, itemsep=3pt, topsep=2pt, parsep=0pt}

\setlength{\parskip}{2pt}
\setlength{\columnsep}{10pt}

\def\CAPETBEGINFILE{\begin{document}}
\def\CAPETENDFILE{\end{document}}
\else
\def\CAPETBEGINFILE{}
\def\CAPETENDFILE{}
\fi

\fancyhead[L]{\textbf{\textcolor{colElec}{T16}} - energies renouvelables : photovoltaeque, hydraulique, eolien}
\fancyhead[R]{\small\textit{CAPET SII IE - Session 2026}}
\fancyfoot[C]{\small\thepage\,/\,\pageref{LastPage}}

\CAPETBEGINFILE

% -----------------------------------------------------------------------
%  BLOC EN-TeTE
% -----------------------------------------------------------------------

\begin{tcolorbox}[
  enhanced, breakable,
  colback=colElec,
  colframe=colElec,
  left=6pt, right=6pt, top=5pt, bottom=5pt,
  before skip=0pt, after skip=4pt
]
  {\color{white}\bfseries\large T16 \quad energies renouvelables : photovoltaeque, hydraulique, eolien}\\[2pt]
  {\color{colElec!30!white}\small
    \textbf{Domaine :} electrotechnique - Systemes energetiques \quad
    \textbf{Presence :} \textbf{7/13} sujets (2013-2026) - theme \emph{frequent}
  }\\[1pt]
  {\color{colElec!20!white}\footnotesize
    Sous-themes : modele panneau PV ($I_{sc}$, $V_{oc}$, MPP) - nombre de panneaux et superficie - MPPT algorithme P\&O - hacheur elevateur - micro-hydraulique (STEP) - eolien (loi de Betz, DFIG)
  }
\end{tcolorbox}

% -----------------------------------------------------------------------
%  1. RAPPELS DE COURS
% -----------------------------------------------------------------------

\begin{tcolorbox}[
  ficheSection=colElec,
  title={1. Rappels de cours},
]

% -- 1.1 Panneau photovoltaeque ----------------------------------------
\textbf{1.1 Modele electrique d'un panneau PV (STC : $1000\,\mathrm{W\,m^{-2}}$, $25^\circ\mathrm{C}$)}

\smallskip
\begin{center}
\renewcommand{\arraystretch}{1.3}
\begin{tabular}{lllll}
\toprule
\textbf{Parametre} & \textbf{Notation} & \textbf{Valeur typ.} & \textbf{En serie} & \textbf{En parallele} \\
\midrule
Puissance crete & $P_{crete}$ & 320\,W & - & - \\
Courant court-circuit & $I_{sc}$ & 9{,}76\,A & $I_{sc}$ inchange & additionne \\
Tension circuit ouvert & $V_{oc}$ & 40{,}96\,V & additionnee & $V_{oc}$ inchange \\
Courant MPP & $I_{MPP}$ & 9{,}34\,A & idem & additionne \\
Tension MPP & $V_{MPP}$ & 34{,}32\,V & additionnee & idem \\
\bottomrule
\end{tabular}
\end{center}
\textbf{Correction en temperature :} $+1^\circ\mathrm{C} \to -0{,}4\,\%$ de $P_{crete}$ (Si-mono). A $50^\circ\mathrm{C}$ (STC + $25^\circ\mathrm{C}$) : $-10\,\%$ de puissance.

Puissance du champ : $P_{champ} = N_{panneaux} \times P_{crete} \times G/G_0$ avec $G_0 = 1000\,\mathrm{W\!\cdot\!m^{-2}}$.

\medskip\hrule\vspace{4pt}

% -- 1.2 MPPT ---------------------------------------------------------
\textbf{1.2 MPPT - Maximum Power Point Tracking}

\smallskip
Commande ajustant le rapport cyclique $\alpha$ du hacheur elevateur pour maintenir le panneau au point de puissance maximale ($V_{MPP}$, $I_{MPP}$).

\textbf{Algorithme perturbation et observation (P\&O) :}
\begin{itemize}[itemsep=1pt]
  \item Imposer $U + \Delta U$ au panneau
  \item Si $\Delta P > 0$ : on se rapproche du MPP - continuer dans le meme sens
  \item Si $\Delta P < 0$ : on s'eloigne - inverser le sens de perturbation
  \item Repeter jusqu'a $\Delta P \approx 0$ (convergence vers le MPP)
\end{itemize}

Rapport cyclique hacheur elevateur : $f_m = V_{PV} / E_b$ (rapport de la tension PV - la tension du bus continu de sortie).

\medskip\hrule\vspace{4pt}

% -- 1.3 Connexion sur onduleur ----------------------------------------
\textbf{1.3 Dimensionnement de la connexion sur onduleur}

\smallskip
\[
  N_{serie/entree} = \frac{V_{MPP\,max\,onduleur}}{V_{MPP\,panneau}}
  \qquad;\qquad
  N_{onduleurs} = \frac{N_{panneaux\,total}}{N_{entrees/onduleur} \times N_{serie/entree}}
\]
Verifier que $V_{MPP}$ et $I_{MPP}$ calcules sont dans les plages minemax de l'onduleur.

\medskip\hrule\vspace{4pt}

% -- 1.4 eolien -------------------------------------------------------
\textbf{1.4 eolien - loi de Betz}

\smallskip
\[
  P_{extrait} = C_p \cdot \frac{1}{2}\,\rho_{air}\,A_{rotor}\,v^3
  \qquad;\qquad
  C_p^{max} = \frac{16}{27} \approx 0{,}593 \quad \text{(limite de Betz)}
\]
$A_{rotor} = \pi R^2$, $v$ = vitesse du vent. Si $v$ double, $P$ est multipliee par \textbf{8}.

\end{tcolorbox}

% -----------------------------------------------------------------------
%  2. FORMULAIRE SYNTHeTIQUE
% -----------------------------------------------------------------------

\begin{tcolorbox}[
  ficheSection=colElec,
  title={2. Formulaire synthetique},
]
\small
\begin{multicols}{2}
\renewcommand{\arraystretch}{1.45}
\begin{tabular}{@{}p{4.5cm}ll@{}}
\toprule
\textbf{Formule} & \textbf{Usage} & \textbf{U.} \\
\midrule
\multicolumn{3}{@{}l}{\footnotesize\textit{Champ PV}} \\
$P_{champ} = N \cdot P_{crete} \cdot G/G_0$ & Puissance & W \\
$N = P_{souhaitee}/P_{crete}$ & Nb panneaux & - \\
$V_{MPP,chaine} = N_s V_{MPP}$ & Tension serie & V \\
\midrule
\multicolumn{3}{@{}l}{\footnotesize\textit{MPPT}} \\
$f_m = V_{PV}/E_b$ & Rapp. cyclique & - \\
P\&O : si $\Delta P>0$ - meme sens & - & - \\
P\&O : si $\Delta P<0$ - inverser & - & - \\
\bottomrule
\end{tabular}

\columnbreak

\renewcommand{\arraystretch}{1.45}
\begin{tabular}{@{}p{4.5cm}ll@{}}
\toprule
\textbf{Formule} & \textbf{Usage} & \textbf{U.} \\
\midrule
\multicolumn{3}{@{}l}{\footnotesize\textit{Hydraulique}} \\
$E_p = \rho V g h$ & energie pot. & J \\
$E_{elec} = E_p\cdot\eta_{turb}$ & Turbinage & Weh \\
$T = V/Q$ & Duree pomp. & h \\
\midrule
\multicolumn{3}{@{}l}{\footnotesize\textit{eolien}} \\
$P = C_p \cdot \frac{1}{2}\rho A v^3$ & Puissance & W \\
$C_p^{max} = 16/27$ & Betz & - \\
\midrule
\multicolumn{3}{@{}l}{\footnotesize\textit{Donnees 2021}} \\
$N = 782$, $S = 1300\,\mathrm{m^2}$ & PV 250\,kW & - \\
$N_s = 26$ panneaux/entree & Onduleur & - \\
\bottomrule
\end{tabular}
\end{multicols}
\normalsize
\end{tcolorbox}

% -----------------------------------------------------------------------
%  3. QUESTIONS TYPES
% -----------------------------------------------------------------------

\begin{tcolorbox}[
  ficheSection=colElec,
  title={3. Questions types - extraites des sujets reels},
]

% -- QT1 -------------------------------------------------------------

% Mini-procedure de resolution type copie concours
\begin{tcolorbox}[questiontype,title={Procedure type copie concours (5 etapes)}]
\begin{enumerate}[leftmargin=1.5em]
  \item \textbf{Qualifier le modele} : hypotheses, domaine de validite, grandeurs d'entree/sortie.
  \item \textbf{Ecrire les lois} : equations de base avec symboles et unites SI.
  \item \textbf{Isoler l'inconnue} : expression litterale, puis application numerique tracee.
  \item \textbf{Controler l'ordre de grandeur} : coherence physique et unite finale.
  \item \textbf{Conclure concours} : phrase courte orientee decision technique.
\end{enumerate}
\end{tcolorbox}

% Visuel de methode : chemin de resolution CAPET
\begin{center}
\begin{tikzpicture}[font=\small, node distance=7mm and 6mm, >=Stealth]
  \node[draw, rounded corners, fill=colElec!8, inner sep=4pt] (n1) {Notions};
  \node[draw, rounded corners, fill=colElec!8, right=of n1, inner sep=4pt] (n2) {Methode};
  \node[draw, rounded corners, fill=colElec!8, right=of n2, inner sep=4pt] (n3) {Application annale};
  \node[draw, rounded corners, fill=colElec!8, right=of n3, inner sep=4pt] (n4) {Verification jury};
  \draw[->, thick, colElec] (n1) -- (n2);
  \draw[->, thick, colElec] (n2) -- (n3);
  \draw[->, thick, colElec] (n3) -- (n4);
\end{tikzpicture}
\end{center}
\begin{tcolorbox}[questiontype,
  title={QT1 - Nombre de panneaux et parametrage cellule \hfill\normalfont\footnotesize E2 disc. 2021 (Q11-Q12)}
]
\textbf{enonce.}
Panneau DMH320M6 : $P_{crete} = 320\,\mathrm{W}$, 60 cellules en serie, $I_{sc} = 9{,}76\,\mathrm{A}$, $V_{oc} = 40{,}96\,\mathrm{V}$, surface 1{,}66\,me. Puissance crete souhaitee : 250\,kW.
(a) Calculer le nombre de panneaux et la superficie. (b) Indiquer $I_{sc,cell}$ et $V_{oc,cell}$ pour une cellule.

\tcbline

\textbf{Corrige.}

\textbf{(a)} $N = 250\,000 / 320 = 782\,\text{panneaux}$ ; superficie $= 782 \times 1{,}66 = \mathbf{1298\,\mathrm{m^2}} \approx 1300\,\mathrm{m^2}$.

\textbf{(b)} Cellules en serie - courant identique - $I_{sc,cell} = I_{sc,panneau} = \mathbf{9{,}76\,\mathrm{A}}$.\\
Tensions s'ajoutent - $V_{oc,cell} = V_{oc,panneau}/60 = 40{,}96/60 = \mathbf{0{,}683\,\mathrm{V}}$.
\end{tcolorbox}

\vspace{4pt}

% -- QT2 -------------------------------------------------------------

\begin{tcolorbox}[questiontype,
  title={QT2 - Dimensionnement entrees onduleur et verification \hfill\normalfont\footnotesize E2 disc. 2021 (Q13), Rapport jury 2021}
]
\textbf{enonce.}
5 onduleurs - 6 entrees MPPT = 30 groupes. Repartir 782 panneaux. Calculer $V_{MPP}$, $I_{MPP}$ et $P_{MPP}$ par entree. Verifier la compatibilite avec l'onduleur.

\tcbline

\textbf{Corrige.}
\[
  N_{s/groupe} = \frac{782}{30} \approx 26\,\text{panneaux par entree}
\]
\[
  V_{MPP,entree} = 26 \times 34{,}32 = \mathbf{892\,\mathrm{V}}
  \qquad;\qquad
  I_{MPP,entree} = 9{,}34\,\mathrm{A}
  \qquad;\qquad
  P_{MPP,entree} = 26 \times 320 = \mathbf{8320\,\mathrm{W}}
\]
\textbf{Verification :} les valeurs calculees ($V_{MPP} = 892\,\mathrm{V}$, $I_{MPP} = 9{,}34\,\mathrm{A}$) doivent etre dans les plages minemax de l'onduleur (verifier sur datasheet constructeur). Le jury insiste sur cette verification explicite.
\end{tcolorbox}

\vspace{4pt}

% -- QT3 -------------------------------------------------------------

\begin{tcolorbox}[questiontype,
  title={QT3 - Algorithme MPPT - perturbation et observation \hfill\normalfont\footnotesize E2 disc. 2025 (Q17-Q18), Rapport jury 2025}
]
\textbf{enonce.}
Sur la courbe $P = f(U)$ d'un champ PV, le point MPP est - $U_{MPP} = 317\,\mathrm{V}$.\\
\textbf{Cas P1 :} $U_1 = 195\,\mathrm{V}$, on impose $\Delta U > 0$. Quel est le signe de $\Delta P$ - Que fait le regulateur MPPT --\\
\textbf{Cas P2 :} $U_2 = 340\,\mathrm{V}$, on impose $\Delta U > 0$. Meme question.

\tcbline

\textbf{Corrige.}

\textbf{Cas P1 ($U_1 = 195\,\mathrm{V} < U_{MPP}$) :} on est - gauche du MPP. En augmentant $U$, on se rapproche du MPP - $\mathbf{\Delta P > 0}$.
Reponse du regulateur : continuer dans le meme sens (augmenter le rapport cyclique du hacheur pour augmenter $U$).

\textbf{Cas P2 ($U_2 = 340\,\mathrm{V} > U_{MPP}$) :} on est - droite du MPP. En augmentant $U$, on s'eloigne - $\mathbf{\Delta P < 0}$.
Reponse du regulateur : inverser le sens (reduire le rapport cyclique pour diminuer $U$).

Le processus converge quand $\Delta P \approx 0$, c'est--dire au voisinage de $U_{MPP} = 317\,\mathrm{V}$.
\end{tcolorbox}

\end{tcolorbox}

% -----------------------------------------------------------------------
%  4. ATTENTES DU JURY
% -----------------------------------------------------------------------

\begin{tcolorbox}[
  jury,
  title={4. Attentes du jury - Erreurs recurrentes et points de vigilance},
]

\begin{enumerate}[label=\textbf{\arabic*.}, leftmargin=*, itemsep=5pt]

  \item \textbf{Correction en temperature obligatoire.}
  Rapport jury 2021 : "les candidats se sont contentes de lire les valeurs sur les courbes sans effectuer la correction en temperature". A $50^\circ\mathrm{C}$, multiplier la puissance STC par $1 - 0{,}004 \times (50 - 25) = 0{,}90$.

  \item \textbf{Verifier la compatibilite avec l'onduleur.}
  Apres avoir calcule $V_{MPP}$ et $I_{MPP}$, verifier explicitement que ces valeurs sont dans les plages admissibles de l'onduleur. Omission signalee comme lacune frequente.

  \item \textbf{Cellules en serie dans un panneau $\neq$ panneaux en serie.}
  Pour une cellule dans un panneau de 60 cellules serie : $I_{sc,cell} = I_{sc,panneau}$ mais $V_{oc,cell} = V_{oc,panneau}/60$. Ne pas confondre avec l'association de panneaux.

  \item \textbf{MPPT : sens de la perturbation.}
  Si $\Delta P > 0$ apres $+\Delta U$ - continuer dans le meme sens. Si $\Delta P < 0$ - inverser. Beaucoup de candidats inversent la logique.

  \item \textbf{Puissance eolienne $\propto v^3$.}
  Si la vitesse du vent double, la puissance est multipliee par $2^3 = 8$, pas par 2. Erreur frequente.

  \item \textbf{Rendement STEP : conclure sur l'interet contextuel.}
  Citer le chiffre numerique (58\,\%) ET la conclusion : "peu performant, mais ecologique, preferable - des batteries supplementaires, exploitable geographiquement".

\end{enumerate}

\medskip
\begin{tcolorbox}[
  colback=colJuryFrame!12!colJury,
  colframe=colJuryFrame,
  left=4pt, right=4pt, top=2pt, bottom=2pt,
  fontupper=\small\itshape
]
Rapport jury 2021 : "Le jury insiste sur la necessite de justifier le choix de l'onduleur en verifiant que les valeurs de tension et de courant calculees sont bien dans les plages admissibles de l'onduleur."\\
Rapport jury 2025 : "MPPT : cette commande permet d'extraire le maximum de puissance. Le rapport cyclique augmente ou diminue en fonction de l'evolution de dP jusqu'a obtenir le point maximum (dP = 0)."
\end{tcolorbox}

\end{tcolorbox}

% -----------------------------------------------------------------------
%  PIED - PReSENCE DANS LES SUJETS
% -----------------------------------------------------------------------

\begin{tcolorbox}[
  enhanced, breakable,
  colback=colGrayLight,
  colframe=colSeparator,
  left=4pt, right=4pt, top=3pt, bottom=3pt,
  before skip=4pt, after skip=0pt,
  fonttitle=\small\bfseries,
  title={Presence dans les sujets analyses (2013-2026)},
]
\begin{center}
\footnotesize\renewcommand{\arraystretch}{1.15}
\begin{tabular}{lccccccccccccc}
\toprule
\textbf{Annee} & 13 & 14 & 15 & 16 & 17 & 18 & 19 & 20 & 21 & 22 & 23 & 24 & 25 \\
\midrule
E1 (commune)   & - & - & - & - & - & - & - & - & - & - & - & - & - \\
E2 discipl. IE & - & - & - & $\checkmark$ & - & - & - & $\checkmark$ & $\checkmark$ & $\checkmark$ & $\checkmark$ & - & $\checkmark$ \\
\bottomrule
\end{tabular}
\end{center}
\smallskip
\footnotesize\textbf{Sous-themes frequents :}
panneau PV - nombre et superficie (4/13) - MPPT algorithme P\&O (3/13) - dimensionnement onduleur (3/13) - micro-hydraulique energie pot. (2/13) - comparaison filieres renouvelables (2/13)

\medskip
\textbf{Ordres de grandeur (reperes rapides).}
\begin{itemize}[leftmargin=1.3em]
  \item Reseau BT FR : tension phase-neutre typique \SI{230}{\volt}, tension phase-phase typique \SI{400}{\volt}, frequence \SI{50}{\hertz}.
  \item Courants d'etude frequents en sujet : de quelques \si{\ampere} a quelques dizaines de \si{\ampere}.
  \item En conversion d'energie, un rendement attendu reste generalement entre 0{,}8 et 0{,}98 selon la technologie.
\end{itemize}

\medskip
\textbf{Synthese express (a refaire sans support).}
\begin{enumerate}[leftmargin=1.4em]
  \item Poser les hypotheses du modele et les conventions de signe.
  \item Ecrire 2 a 4 relations cles avec unites.
  \item Realiser un calcul numerique complet et verifier l'ordre de grandeur.
  \item Justifier le resultat par un argument technique de niveau CAPET.
\end{enumerate}

\end{tcolorbox}

\CAPETENDFILE









